SRQ1 & Explicit belief over a hidden incident state is used by a minority of studies & \cite{hammar2026belief,hammar2025,koupaei2026,kovtun2026} & The representation the project needs exists and has been evaluated, mostly on security benchmarks \\
 & Ambiguity is produced mainly by noisy, missing or overlapping alerts in simulators & \cite{koupaei2026,contractor2025,sahin2026,wu2026,peralta2025,aenugu2026,ravi2025} & Few studies vary ambiguity as an experimental factor; most fix it at one level \\
 & Many studies treat the incident state as a fully observed feature vector & \cite{saeed2025,noorulhassan2026,alturkistani2022,bhosale2026,raj2026,nimmathi2026,yuan2026,babar2026,hughes2022,jaber2025,singh2026,ghanem2026,dunsin2025,wang2024mrca,ding2023} & Ambiguity is often assumed away rather than modelled \\
 & Evaluation environments are mostly security benchmarks, custom simulators and repurposed IDS datasets & \cite{hammar2026belief,hammar2025,castro2025,contractor2025,nyberg2024,rojasrodriguez2026,okafor2026,kovtun2026,sahin2026,peralta2025,wang2023hrlhf,ding2023} & No shared environment exists for IT service incident triage \\
\midrule
SRQ2 & Information gathering is an explicit action in about half the studies & \cite{okafor2026,hammar2026belief,rojasrodriguez2026,sahin2026,tong2020,castro2025,contractor2025,ghanem2026,jaber2025,wang2024mrca,hammar2025,dunsin2025,wu2026,kovtun2026,wang2023hrlhf,nyberg2024,ding2023} & Investigation is treated as a decision with a cost, as the project assumes \\
 & Escalation to a human is rarely a learned action & \cite{bhosale2026,alturkistani2022,sahin2026,kovtun2026,okafor2026,singh2026,koupaei2026} & Where escalation exists it is mostly in security alert triage, and S038 hands off by rule rather than by learning \\
 & No study learns one policy over investigating, remediating and escalating for IT service incidents & \cite{sahin2026,okafor2026,kovtun2026} & The three-way decision of the project is not covered; the nearest work is security alert triage \\
 & When investigation is available it reduces errors or effort & \cite{wu2026,kovtun2026,ghanem2026,wang2023hrlhf} & Supports giving the agent a costly investigate action \\
\midrule
SRQ3 & Hard runtime filters and shields are proposed but rarely measured & \cite{ravi2025,singh2026,koupaei2026,elicha2026,okafor2026} & Only S118 and S038 report how often rules intervened or were complied with \\
 & Action masking and expected-cost constraints are the common lightweight alternatives & \cite{contractor2025,rojasrodriguez2026,wang2024mrca,ghanem2026,wu2026,ding2023,kovtun2026,tong2020,varsha2025,sahin2026} & Masking prevents prohibited proposals rather than reacting to them \\
 & What happens to a blocked action is reported in only three studies & \cite{ravi2025,singh2026,sahin2026} & The project's rule of one re-proposal then escalation has no direct precedent \\
 & Rule enforcement changes behaviour and has a cost & \cite{singh2026,ravi2025} & Compliance has to be measured alongside decision quality and cost \\
\midrule
SRQ4 & Decision quality and time dominate the metrics; escalation and compliance are rarely measured & All 34 (Table~\ref{tab:quality}) & Only one study measures all four groups the project plans, and it concerns SOC alert escalation \\
 & A static or rule-based comparison is common but not universal & \cite{peralta2025,hammar2026belief,singh2026,babar2026,rojasrodriguez2026,koupaei2026,sahin2026,tong2020,nimmathi2026,castro2025,ghanem2026,varsha2025,wang2024mrca,aenugu2026,bhosale2026,dunsin2025,wu2026,wang2023hrlhf,yuan2026,ding2023} & A rule-based baseline is standard and the project should keep one \\
 & Reporting weaknesses limit the comparisons & \cite{okafor2026,ravi2025,koupaei2026,bhosale2026,varsha2025,jaber2025,alturkistani2022,sahin2026} & Metrics such as MCC, per-class recall and seed variance are needed \\
\midrule
SRQ5 & Robustness is tested through model error, change, unseen attackers and weaker detection & \cite{hammar2025,hammar2026belief,ravi2025,nyberg2024,contractor2025,sahin2026,peralta2025,wu2026,aenugu2026,castro2025,kovtun2026} & Robustness is tested in about a third of studies, mostly for decision quality \\
 & Methods that maintain a belief or adapt their model degrade least & \cite{hammar2025,hammar2026belief} & Supports explicit belief tracking in the project agent \\
 & As ambiguity rises, deferral to people rises; how the full decision mix shifts is unmeasured & \cite{sahin2026,kovtun2026,contractor2025} & No study reports the shares of investigating, acting and escalating across ambiguity levels \\
